\chapter{Homotopies of Maps}
\label{ch:viii01-homotopies}

Algebraic topology begins by weakening the notion of equality of maps.  Two
continuous maps need not agree pointwise to carry the same large-scale
topological information.  The basic relation is \emph{homotopy}: one map may be
continuously deformed into the other while the domain and target remain fixed.
This chapter develops the definition, the elementary calculus of homotopies,
relative and based variants, and the first explicit deformation formulas.

\section{The unit interval as a deformation parameter}

Write
\[
I=[0,1].
\]
A point of \(X\times I\) is written \((x,t)\), with \(t\) interpreted as time.
At \(t=0\) we see the initial map and at \(t=1\) the final map.

\begin{definition}[Homotopy]
\label{def:viii01-homotopy}
Let \(f,g:X\to Y\) be continuous maps.  A \emph{homotopy from \(f\) to \(g\)}
is a continuous map
\[
H:X\times I\longrightarrow Y
\]
such that
\[
H(x,0)=f(x),
\qquad
H(x,1)=g(x)
\]
for every \(x\in X\).  We write
\[
f\simeq g.
\]
\end{definition}

For each fixed \(t\), the map
\[
H_t:X\to Y,
\qquad
H_t(x)=H(x,t),
\]
is a time slice of the deformation.

\begin{warning}[Pointwise paths are not enough]
\label{warn:viii01-pointwise}
For every \(x\), the curve \(t\mapsto H(x,t)\) is a path in \(Y\), but choosing
an arbitrary path from \(f(x)\) to \(g(x)\) for each \(x\) does not automatically
produce a homotopy.  The dependence on \(x\) and \(t\) must be jointly continuous.
\end{warning}

\section{First explicit homotopies}

\begin{example}[Straight-line homotopy]
\label{ex:viii01-straight}
If \(Y\) is a convex subset of a real topological vector space and
\(f,g:X\to Y\) are continuous, then
\[
H(x,t)=(1-t)f(x)+tg(x)
\]
is a homotopy from \(f\) to \(g\).
\end{example}

\begin{example}[Maps into an interval]
\label{ex:viii01-interval-target}
Any two continuous maps \(f,g:X\to [a,b]\) are homotopic by the straight-line
formula.  Thus the interval has no obstruction to deforming one incoming map
into another.
\end{example}

\begin{example}[Rotation of the circle]
\label{ex:viii01-circle-rotation}
For a fixed angle \(\alpha\), the identity map on \(S^1\subset\mathbb C\) is
homotopic to the rotation \(R_\alpha(z)=e^{i\alpha}z\) via
\[
H(z,t)=e^{it\alpha}z.
\]
\end{example}

\section{Homotopy is an equivalence relation}

\begin{theorem}[Equivalence relation]
\label{thm:viii01-equivalence}
Homotopy is reflexive, symmetric, and transitive on the set of continuous maps
\(X\to Y\).
\end{theorem}

\begin{proof}
Reflexivity is given by the constant homotopy \(H(x,t)=f(x)\).
If \(H\) runs from \(f\) to \(g\), then
\[
\overline H(x,t)=H(x,1-t)
\]
runs from \(g\) to \(f\).  For transitivity, suppose \(H:f\simeq g\) and
\(K:g\simeq h\).  Define
\[
L(x,t)=
\begin{cases}
H(x,2t),&0\le t\le \frac12,\\[2mm]
K(x,2t-1),&\frac12\le t\le1.
\end{cases}
\]
The two formulas agree at \(t=\frac12\), where both equal \(g(x)\).  The gluing
lemma gives continuity.
\end{proof}

\begin{remark}[Homotopy classes]
\label{rem:viii01-classes}
The equivalence class of \(f:X\to Y\) is denoted \([f]\).  The set of homotopy
classes is often written
\[
[X,Y].
\]
At this stage it is only a set; later chapters supply algebraic structure in
important cases.
\end{remark}

\section{Compatibility with composition}

\begin{proposition}[Postcomposition]
\label{prop:viii01-postcomposition}
If \(f\simeq g:X\to Y\) and \(u:Y\to Z\) is continuous, then
\[
u\circ f\simeq u\circ g.
\]
\end{proposition}

\begin{proof}
If \(H\) is a homotopy from \(f\) to \(g\), use
\[
(x,t)\longmapsto u(H(x,t)).
\]
\end{proof}

\begin{proposition}[Precomposition]
\label{prop:viii01-precomposition}
If \(f\simeq g:X\to Y\) and \(v:W\to X\) is continuous, then
\[
f\circ v\simeq g\circ v.
\]
\end{proposition}

\begin{proof}
Use
\[
(w,t)\longmapsto H(v(w),t).
\]
\end{proof}

\begin{corollary}[Two-sided compatibility]
\label{cor:viii01-two-sided}
If \(f_0\simeq f_1:X\to Y\) and \(g_0\simeq g_1:Y\to Z\), then
\[
g_0\circ f_0\simeq g_1\circ f_1.
\]
\end{corollary}

\section{Null-homotopic maps}

\begin{definition}[Null-homotopic map]
\label{def:viii01-null}
A map \(f:X\to Y\) is \emph{null-homotopic} if it is homotopic to a constant
map.
\end{definition}

\begin{example}[Maps into a convex set]
\label{ex:viii01-convex-null}
Every map into a convex subset \(C\subset\mathbb R^n\) is null-homotopic:
choose \(c_0\in C\) and use
\[
H(x,t)=(1-t)f(x)+tc_0.
\]
\end{example}

\begin{proposition}[Factorization through a point]
\label{prop:viii01-point-factor}
A map is null-homotopic precisely when it is homotopic to a map factoring as
\[
X\longrightarrow \{*\}\longrightarrow Y.
\]
\end{proposition}

\section{Homotopies relative to a subspace}

Often a deformation must leave part of the domain fixed.

\begin{definition}[Homotopy relative to \(A\)]
\label{def:viii01-relative}
Let \(A\subset X\).  A homotopy \(H:f\simeq g\) is a \emph{homotopy relative to
\(A\)}, written
\[
f\simeq g\pmod A,
\]
if
\[
H(a,t)=f(a)=g(a)
\]
for all \(a\in A\) and \(t\in I\).
\end{definition}

\begin{remark}
A relative homotopy can exist only if \(f|_A=g|_A\).  The condition says more:
the common restriction is fixed throughout the deformation.
\end{remark}

\begin{example}[Fixing endpoints]
\label{ex:viii01-fixed-endpoints}
For paths \(\alpha,\beta:I\to Y\) with common endpoints, a homotopy relative to
\(\{0,1\}\) is exactly an endpoint-fixing deformation of one path into the other.
This is the relation used later in the fundamental group.
\end{example}

\section{Based homotopy}

Let \((X,x_0)\) and \((Y,y_0)\) be pointed spaces.

\begin{definition}[Based map and based homotopy]
\label{def:viii01-based}
A map \(f:(X,x_0)\to(Y,y_0)\) is based if \(f(x_0)=y_0\).  A homotopy through
based maps is a homotopy relative to \(\{x_0\}\):
\[
H(x_0,t)=y_0
\qquad\text{for all }t.
\]
\end{definition}

\begin{warning}[Free versus based homotopy]
\label{warn:viii01-free-based}
Two based maps can be freely homotopic without being based-homotopic.  The
distinction becomes essential when loops and fundamental groups enter.
\end{warning}

\section{Reparameterizing a homotopy}

\begin{lemma}[Monotone reparameterization]
\label{lem:viii01-reparam}
If \(H:f\simeq g\) and \(\rho:I\to I\) is continuous with
\(\rho(0)=0\) and \(\rho(1)=1\), then
\[
K(x,t)=H(x,\rho(t))
\]
is another homotopy from \(f\) to \(g\).
\end{lemma}

This permits slowing, speeding, or pausing a deformation without changing its
homotopy class.

\section{Homotopy through a subspace}

\begin{proposition}[Image control]
\label{prop:viii01-image-control}
Suppose \(A\subset Y\), and a homotopy \(H:X\times I\to Y\) satisfies
\(H(X\times I)\subset A\).  Then \(H\), regarded as a map into \(A\), is a
homotopy in the subspace \(A\).
\end{proposition}

The proposition is elementary, but important: whether a deformation is allowed
depends on the target in which the deformation is required to live.

\section{A punctured-plane obstruction preview}

The straight-line formula can fail when the target is not convex.

\begin{example}[Why the ambient target matters]
\label{ex:viii01-punctured}
In \(\mathbb R^2\setminus\{0\}\), the formula
\[
H(x,t)=(1-t)x+t(-x)
\]
passes through \(0\) at \(t=\frac12\).  Therefore it is not a homotopy in the
punctured plane.  Later invariants will explain why certain maps cannot be
deformed while avoiding the missing point.
\end{example}

\section{Cylinder viewpoint}

A homotopy \(H:X\times I\to Y\) is simply a map out of the cylinder on \(X\),
with prescribed values on the two ends
\[
X\times\{0\},
\qquad
X\times\{1\}.
\]
This point of view anticipates mapping cylinders, cofibrations, and relative
homotopy constructions.

\begin{proposition}[Restriction to time intervals]
\label{prop:viii01-restrict-time}
If \(H:f\simeq g\) and \(0\le a<b\le1\), then after affine reparameterization
the restriction of \(H\) to \(X\times[a,b]\) is a homotopy from \(H_a\) to
\(H_b\).
\end{proposition}

\section{Products of homotopies}

\begin{proposition}[Product homotopy]
\label{prop:viii01-product}
If \(f_0\simeq f_1:X\to Y\) and \(g_0\simeq g_1:W\to Z\), then
\[
f_0\times g_0\simeq f_1\times g_1:X\times W\to Y\times Z.
\]
\end{proposition}

\begin{proof}
For homotopies \(H\) and \(K\), set
\[
L((x,w),t)=\bigl(H(x,t),K(w,t)\bigr).
\]
\end{proof}

\section{What a homotopy does not claim}

A homotopy does not assert that any \(H_t\) is injective, surjective, an
embedding, a homeomorphism, or a diffeomorphism.  It is a deformation in the
space of continuous maps.  If every time slice is required to be an embedding
or homeomorphism, one obtains stronger notions such as isotopy.

\begin{warning}[Homotopy is weaker than isotopy]
\label{warn:viii01-isotopy}
Do not infer geometric rigidity from a homotopy.  Intermediate maps may fold,
collapse, or identify points unless additional hypotheses prohibit this.
\end{warning}

\section{Exercises}

\begin{exercise}\label{exr:viii01-01}
Write the defining endpoint equations for a homotopy \(H:f\simeq g\).
\end{exercise}
\begin{hint}Evaluate at \(t=0\) and \(t=1\).\end{hint}
\begin{solution}
They are \(H(x,0)=f(x)\) and \(H(x,1)=g(x)\) for every \(x\in X\).
\end{solution}

\begin{exercise}\label{exr:viii01-02}
Show that every map is homotopic to itself.
\end{exercise}
\begin{hint}Do not move the map in time.\end{hint}
\begin{solution}
Use \(H(x,t)=f(x)\).
\end{solution}

\begin{exercise}\label{exr:viii01-03}
Reverse a given homotopy \(H:f\simeq g\).
\end{exercise}
\begin{hint}Reverse the interval parameter.\end{hint}
\begin{solution}
Set \(\overline H(x,t)=H(x,1-t)\).
\end{solution}

\begin{exercise}\label{exr:viii01-04}
Concatenate homotopies \(H:f\simeq g\) and \(K:g\simeq h\).
\end{exercise}
\begin{hint}Give each homotopy half the time interval.\end{hint}
\begin{solution}
Use \(H(x,2t)\) for \(t\le1/2\) and \(K(x,2t-1)\) for \(t\ge1/2\).
\end{solution}

\begin{exercise}\label{exr:viii01-05}
Why is the concatenated homotopy continuous at \(t=1/2\)?
\end{exercise}
\begin{hint}Compare the two values there.\end{hint}
\begin{solution}
Both pieces equal \(g(x)\) at the joining time, so the gluing lemma applies.
\end{solution}

\begin{exercise}\label{exr:viii01-06}
Give a homotopy from \(f(x)=x\) to \(g(x)=0\) on \(\mathbb R^n\).
\end{exercise}
\begin{hint}Use scalar multiplication.\end{hint}
\begin{solution}
\(H(x,t)=(1-t)x\).
\end{solution}

\begin{exercise}\label{exr:viii01-07}
Show that two maps \(X\to[a,b]\) are homotopic.
\end{exercise}
\begin{hint}The interval is convex.\end{hint}
\begin{solution}
Use \(H(x,t)=(1-t)f(x)+tg(x)\).
\end{solution}

\begin{exercise}\label{exr:viii01-08}
Construct a homotopy from the identity of \(S^1\) to rotation by angle \(\pi/3\).
\end{exercise}
\begin{hint}Rotate gradually.\end{hint}
\begin{solution}
\(H(z,t)=e^{it\pi/3}z\).
\end{solution}

\begin{exercise}\label{exr:viii01-09}
If \(f\simeq g\), prove \(u\circ f\simeq u\circ g\).
\end{exercise}
\begin{hint}Postcompose the homotopy.\end{hint}
\begin{solution}
For \(H:f\simeq g\), use \(K(x,t)=u(H(x,t))\).
\end{solution}

\begin{exercise}\label{exr:viii01-10}
If \(f\simeq g\), prove \(f\circ v\simeq g\circ v\).
\end{exercise}
\begin{hint}Precompose the homotopy.\end{hint}
\begin{solution}
Use \(K(w,t)=H(v(w),t)\).
\end{solution}

\begin{exercise}\label{exr:viii01-11}
Define a null-homotopic map.
\end{exercise}
\begin{hint}Compare with a constant map.\end{hint}
\begin{solution}
It is a map homotopic to a constant map.
\end{solution}

\begin{exercise}\label{exr:viii01-12}
Show every map into a convex subset of \(\mathbb R^n\) is null-homotopic.
\end{exercise}
\begin{hint}Choose a point in the convex set.\end{hint}
\begin{solution}
For \(c_0\) in the target, \(H(x,t)=(1-t)f(x)+tc_0\) remains in the convex set.
\end{solution}

\begin{exercise}\label{exr:viii01-13}
State what \(f\simeq g\pmod A\) means.
\end{exercise}
\begin{hint}The subspace must not move.\end{hint}
\begin{solution}
It means there is a homotopy \(H:f\simeq g\) with
\(H(a,t)=f(a)=g(a)\) for all \(a\in A\).
\end{solution}

\begin{exercise}\label{exr:viii01-14}
What preliminary condition on \(f\) and \(g\) is necessary for a homotopy rel \(A\)?
\end{exercise}
\begin{hint}Set \(t=0,1\) on \(A\).\end{hint}
\begin{solution}
They must agree on \(A\): \(f|_A=g|_A\).
\end{solution}

\begin{exercise}\label{exr:viii01-15}
Interpret homotopy rel \(\{0,1\}\) for paths.
\end{exercise}
\begin{hint}The endpoints are the chosen subspace.\end{hint}
\begin{solution}
It is a deformation of paths that keeps both endpoints fixed at all times.
\end{solution}

\begin{exercise}\label{exr:viii01-16}
What is a based homotopy?
\end{exercise}
\begin{hint}Fix the basepoint throughout time.\end{hint}
\begin{solution}
It is a homotopy \(H\) satisfying \(H(x_0,t)=y_0\) for every \(t\).
\end{solution}

\begin{exercise}\label{exr:viii01-17}
Why can an arbitrary collection of paths from \(f(x)\) to \(g(x)\) fail to be a homotopy?
\end{exercise}
\begin{hint}Continuity has two variables.\end{hint}
\begin{solution}
The chosen paths may fail to depend continuously on \(x\), so the resulting
map \(X\times I\to Y\) need not be continuous.
\end{solution}

\begin{exercise}\label{exr:viii01-18}
Reparameterize a homotopy so that it pauses at \(t=1/2\).
\end{exercise}
\begin{hint}Choose a continuous \(\rho:I\to I\) that is constant on a small interval.\end{hint}
\begin{solution}
Choose any continuous endpoint-preserving \(\rho\) that is constant near
\(1/2\), and use \(K(x,t)=H(x,\rho(t))\).
\end{solution}

\begin{exercise}\label{exr:viii01-19}
Show that the product of two null-homotopic maps is null-homotopic.
\end{exercise}
\begin{hint}Take the product of their homotopies.\end{hint}
\begin{solution}
If \(H\) and \(K\) contract the two maps to constants, then
\((H,K)\) contracts their product to the product constant.
\end{solution}

\begin{exercise}\label{exr:viii01-20}
Why does the straight-line deformation from \(x\) to \(-x\) fail in
\(\mathbb R^2\setminus\{0\}\)?
\end{exercise}
\begin{hint}Evaluate at \(t=1/2\).\end{hint}
\begin{solution}
It equals \(0\) halfway through, which is not in the target.
\end{solution}

\begin{exercise}\label{exr:viii01-21}
If \(H:X\times I\to A\subset Y\), in which two targets can the same formula be viewed?
\end{exercise}
\begin{hint}Use the inclusion \(A\hookrightarrow Y\).\end{hint}
\begin{solution}
It is a homotopy into \(A\), and after composing with the inclusion it is also
a homotopy into \(Y\).
\end{solution}

\begin{exercise}\label{exr:viii01-22}
Does a homotopy require every time slice to be injective?
\end{exercise}
\begin{hint}Return to the definition.\end{hint}
\begin{solution}
No.  Only continuity and the endpoint conditions are required.
\end{solution}

\begin{exercise}\label{exr:viii01-23}
Does a homotopy between homeomorphisms have to pass through homeomorphisms?
\end{exercise}
\begin{hint}Compare homotopy with isotopy.\end{hint}
\begin{solution}
No.  That stronger requirement belongs to isotopy.
\end{solution}

\begin{exercise}\label{exr:viii01-24}
Explain why composition descends to homotopy classes.
\end{exercise}
\begin{hint}Use pre- and postcomposition compatibility.\end{hint}
\begin{solution}
Replacing either factor by a homotopic map changes the composite only by a
homotopy, so the homotopy class of a composite depends only on the classes of
its factors.
\end{solution}

\section{Solved problem dossiers}

\begin{problem}[Piecewise concatenation]\label{prob:viii01-01}
Verify directly that concatenating two relative homotopies rel \(A\) again gives
a homotopy rel \(A\).
\end{problem}
\begin{solution}
Use the usual two-piece concatenation.  On \(A\), both original homotopies are
constant at the common restriction \(f|_A=g|_A=h|_A\), so both pieces of the
concatenation are constant there as well.
\end{solution}

\begin{problem}[Radial contraction]\label{prob:viii01-02}
Let \(D^n=\{x:\|x\|\le1\}\).  Construct a homotopy from the identity of \(D^n\)
to the constant map at \(0\).
\end{problem}
\begin{solution}
Set \(H(x,t)=(1-t)x\).  Since
\(\|H(x,t)\|=(1-t)\|x\|\le1\), the image remains in \(D^n\).
\end{solution}

\begin{problem}[Affine homotopy in a convex target]\label{prob:viii01-03}
Prove that the straight-line formula is jointly continuous.
\end{problem}
\begin{solution}
The maps \((x,t)\mapsto (1-t)f(x)\) and \((x,t)\mapsto tg(x)\) are continuous
by continuity of scalar multiplication and of \(f,g\); their sum is continuous.
Convexity ensures the value remains in the target.
\end{solution}

\begin{problem}[Homotopy of constant maps]\label{prob:viii01-04}
If \(Y\) is path connected, prove that any two constant maps \(X\to Y\) are
homotopic.
\end{problem}
\begin{solution}
If the constants are \(y_0,y_1\), choose a path \(\gamma:I\to Y\) from
\(y_0\) to \(y_1\) and define \(H(x,t)=\gamma(t)\).
\end{solution}

\begin{problem}[Path components detected by maps from a point]\label{prob:viii01-05}
Describe \([\{*\},Y]\).
\end{problem}
\begin{solution}
A map from a point chooses a point of \(Y\).  Two such maps are homotopic
exactly when the chosen points are joined by a path.  Hence
\([\{*\},Y]\) is the set of path components of \(Y\).
\end{solution}

\begin{problem}[Restriction of a relative homotopy]\label{prob:viii01-06}
If \(H:f\simeq g\pmod A\) and \(B\subset A\), show \(H\) is a homotopy rel \(B\).
\end{problem}
\begin{solution}
The defining fixed-point equation on all of \(A\) in particular holds on its
subset \(B\).
\end{solution}

\begin{problem}[Composition of based homotopies]\label{prob:viii01-07}
Show that composing based homotopies produces a based homotopy.
\end{problem}
\begin{solution}
If \(H(x_0,t)=y_0\) and \(K(y_0,t)=z_0\), the standard two-sided compatibility
construction keeps the distinguished point at \(z_0\) throughout.
\end{solution}

\begin{problem}[A stationary first half]\label{prob:viii01-08}
Given \(H:f\simeq g\), construct a homotopy that stays equal to \(f\) until
time \(1/2\) and then performs \(H\).
\end{problem}
\begin{solution}
Use
\[
K(x,t)=
\begin{cases}
f(x),&0\le t\le1/2,\\
H(x,2t-1),&1/2\le t\le1.
\end{cases}
\]
The formulas agree at \(t=1/2\).
\end{solution}

\begin{problem}[Product deformation]\label{prob:viii01-09}
If \(X\) and \(W\) each admit a contraction to a chosen point, write the
contraction of \(X\times W\).
\end{problem}
\begin{solution}
For contractions \(H\) and \(K\), use
\[
L((x,w),t)=(H(x,t),K(w,t)).
\]
\end{solution}

\begin{problem}[Target inclusion]\label{prob:viii01-10}
Let \(A\subset Y\) and suppose \(f,g:X\to A\) are homotopic as maps into \(A\).
Prove they are homotopic as maps into \(Y\).
\end{problem}
\begin{solution}
Compose the homotopy \(X\times I\to A\) with the inclusion \(A\hookrightarrow Y\).
\end{solution}

\begin{problem}[Two-stage deformation]\label{prob:viii01-11}
Suppose \(f\simeq c\) and \(c\simeq g\), where \(c\) is constant.  Prove
\(f\simeq g\).
\end{problem}
\begin{solution}
This is transitivity: concatenate the two homotopies.
\end{solution}

\begin{problem}[Bridge to VIII/02]\label{prob:viii01-12}
Explain why a homotopy \(f\simeq g\) alone does not say that the spaces \(X\)
and \(Y\) have the same homotopy type.
\end{problem}
\begin{solution}
The relation compares two maps with the same domain and target.  To compare
spaces, one needs maps in both directions whose composites are homotopic to the
appropriate identity maps.  That is the homotopy-equivalence notion developed
in VIII/02.
\end{solution}
